\subsection{The fixed-exit continuation}
\label{sec:fixed-exit-comparison}

For nonzero limiting action $\Theta=\beta\log V$, orbits of
\eqref{fx:pd}--\eqref{fx:graph} leave the small-height region before
the negative saddle. At the fixed exit $u_c\in(-1,u_-)$ the height
$W=w(u_c)$ provides a regular coordinate. Shooting backwards from this
section constructs the scalar inverse phase and its $C^2$ limit without
using the positive return-section map.

\subsection{The action interval and the exit height}

For \(J(u)=\exp\int_{u_+}^u L/h\), put
\begin{equation}\label{fx:rho-gamma}
 \rho=\frac{h(u_-)J(u)}{J(u_-)h(u)},\qquad
 \Gamma_\beta=-\beta\log\rho,\qquad H_\beta=h/\beta .
\end{equation}
The exact partial fractions are
\begin{equation}\label{fx:partial-fractions}
 \frac Lh=-\frac1{\beta(u+1)}
       +\frac1{\chi(a-u)}
       -\frac{1+1/\beta+1/\chi}{c-u}.
\end{equation}
Thus
\begin{equation}\label{fx:gamma-derivative}
 \Gamma_\beta'=-L/H_\beta+\beta H_\beta'/H_\beta,\qquad
 H_\beta=\frac{b_-(u+1)(a-u)(c-u)}{(1+a)(1+c)} .
\end{equation}
On negative compact intervals this tends in \(C^3\) to
\begin{equation}\label{fx:limit-functions}
 H_{0,c}=\frac{(-u)(1+u)(c-u)}{1+c},\qquad
 \Gamma_c(u)=\log\frac{(1+u)(c-u_-)}{(1+u_-)(c-u)} .
\end{equation}
Indeed, \(\Gamma_c'=-u/H_{0,c}\), and both actions vanish at \(u_-\).
The limit action labels the point at which an orbit leaves the
small-height region. From that point to the fixed exit, the limiting
graph equation is \(w_u=u\). These two limits will parametrize the
inverse transition.

At the shrinking positive cut, integration of
\eqref{fx:partial-fractions} gives
\[
 \beta\log J(u_-)=
 \log\frac{(1+u_+)(c-u_-)}{(1+u_-)(c-u_+)}
 +q\log\frac{(a-u_+)(c-u_-)}{(a-u_-)(c-u_+)}
 +\beta\log\frac{c-u_-}{c-u_+}.
\]
By \eqref{fx:log-smallness},
\begin{equation}\label{fx:gamma-endpoint}
 \Gamma_\beta(u_+)\longrightarrow
 \Gamma_{-,c}:=\Gamma_c(0)=
             \log\frac{c-u_-}{c(1+u_-)}>0 .
\end{equation}
Also \((L-h')'=1+o(1)>0\); hence
\(\Gamma_\beta'=-\beta(L-h')/h\) has at most one zero, a maximum.

Choose \(0<\Delta<\frac14\inf_c\Gamma_{-,c}\), and set
\[
 r_1=\Gamma_c^{-1}(\Gamma_{-,c}-\Delta),\qquad
 r_1^\sharp=\Gamma_c^{-1}(\Gamma_{-,c}-3\Delta/4).
\]
Choose fixed cuts \(u_c,u_R\) and \(\epsilon>0\) with
\[
 -1<u_c<u_--2\epsilon,\qquad
 r_1^\sharp(c)<u_R<0,\qquad
 \Gamma_c(u_R)>\Gamma_{-,c}-\Delta/2 .
\]
These choices are uniform in the compact \(c\)-family. Define
\begin{equation}\label{fx:W-limit}
 \mathcal W(r)=\frac{u_c^2-r^2}{2},\quad
 \mathcal J_c=\mathcal W([u_-,r_1]),\quad
 \mathcal J_c^\sharp=\mathcal W([u_--\epsilon,r_1^\sharp]).
\end{equation}
They are positive compact intervals with a uniform strict inclusion.

\subsection{Convergence of the inverse transition}

We first work only with the scalar equation \eqref{fx:graph}, without
assuming any positive return-section map. Let \(\mathcal K_c\) be a
compact interval strictly inside \((0,u_c^2/2)\), with uniform endpoint
margins, and prescribe \(w(u_c;W)=W\) for \(W\in\mathcal K_c\).
Write
\[
 r(W)=-\sqrt{u_c^2-2W},\qquad
 d_{\mathcal K}=\inf_{c,W\in\mathcal K_c}
       \{\Gamma_{-,c}-\Gamma_c(r(W))\}>0.
\]
Increase \(u_R<0\) towards zero if necessary, so that it lies beyond
all \(r(W)\) and
\(\Gamma_c(u_R)>\Gamma_{-,c}-d_{\mathcal K}/4\).
We follow the backward shot through its positive-height part, the thin
transition layer, and the remaining small-height part. The next three
first-exit arguments continue it to \(u_+\) and estimate its phase.
They apply, in particular, to
\(\mathcal J_c^\sharp\) and to any fixed larger chord still inside
\((0,u_c^2/2)\). Define
\begin{equation}\label{fx:extracted-phase}
 \phi_\beta(u;W)=\beta\log\frac{w(u;W)}{h(u)}+\Gamma_\beta(u).
\end{equation}
Its value at \(u_+\) gives the exact inverse phase
\begin{equation}\label{fx:inverse-phase}
 \Psi_\theta(W)=\phi_\beta(u_+;W)-\beta\log\lambda
                 =\beta\log V=:\Theta .
\end{equation}
The differential equation for this phase is
\begin{equation}\label{fx:phase-equation}
 \phi_\beta'=\frac{r(Eh-KL)}{H_\beta(1+Kr)},\qquad r=w/h .
\end{equation}
The phase changes by only a small amount after the orbit has entered
the small-height part. This will allow the leaving point of the limiting
graph to determine the inverse phase.

\begin{lemma}\label{fx:lem-fast-inverse}
For fixed \(\omega_0>0\), the graph meets \(w=\beta\omega_0\)
at a unique point \(\sigma_\beta(W)<u_R\), and, uniformly on
\(\mathcal K_c\),
\[
 \|\sigma_\beta-\mathcal W^{-1}\|_{C^2}
     \le C(\omega_{\rm d}+\beta|\log\beta|)\longrightarrow0 .
\]
Here \(\omega_{\rm d}\) bounds in \(C^3\) the differences
\(H_\beta-H_{0,c},K-1,L-u,E\) on the fixed negative rectangle.
\end{lemma}
\begin{proof}
On that rectangle \(L+Ew\le-\ell<0\).
The inverse graph, \(U_\beta(W;W)=u_c\), satisfies
\[
 \partial_w U_\beta
       =\frac K{L+Ew}+\frac{\beta H_\beta}{w(L+Ew)} .
\]
Between \(\beta\omega_0\) and the fixed upper height, its
right side differs from \(1/u\) in \(L^1_w\) by at most
\(C\omega_{\rm d}+C\int_{\beta\omega_0}^{W_{\max}}\beta\,dw/w\).
Gronwall's inequality and the first two triangular variational
equations give the asserted \(C^2\) bound.
To justify existence, apply this estimate only up to a first exit
from a fixed rectangle about the limiting graphs. Their endpoints
\(r(W)<u_R\) stay strictly inside,
so an exit is impossible. Negative nonzero speed at the selected
height gives uniqueness and transversality.
\end{proof}

\begin{lemma}\label{fx:lem-corner}
Uniformly on \(\mathcal K_c\),
\[
 \|\phi_\beta(u_R;W)-\Gamma_\beta(\sigma_\beta(W))\|_{C^2}
       \le C\beta .
\]
\end{lemma}
\begin{proof}
On this fixed negative rectangle,
\(g_\beta=-(L-h')/H_\beta\ge g_0>0\), and the exact normal equation is
\[
 \beta r'=\frac{r\{-g_\beta+\beta r(E-KH_\beta'/H_\beta)\}}{1+Kr}.
\]
Set \(\zeta=(u-\sigma)/\beta\), so \(r_\zeta=f_\beta\),
\(r(0;\sigma)=\omega_0/H_\beta(\sigma)\). On the fixed bounded
\(r\)-interval, \(f_r\le-a_0<0\),
\(|f_u|+|f_{uu}|\le Cr\), and \(|f_{ur}|+|f_{rr}|\le C\).
The triangular variation argument of
Lemma~\ref{lem:scalar-normalization}, with these decaying parameter
forcings, gives
\[
 |\partial_\sigma^j r|\le C_j(1+\zeta^j)e^{-c\zeta},\qquad j=0,1,2.
\]
A first-exit argument retains the rectangle up to \(u_R\).
For \(G(\sigma,\zeta)=F_\beta(\sigma+\beta\zeta,r)\), where
\(F_\beta\) is the right side of \eqref{fx:phase-equation}, put
\(T=(u_R-\sigma)/\beta\), \(J=\beta\int_0^T G\,d\zeta\).
The moving endpoint must be retained:
\[
 J_\sigma=\beta\int_0^T G_\sigma\,d\zeta-G(\sigma,T),\qquad
 J_{\sigma\sigma}=\beta\int_0^T G_{\sigma\sigma}\,d\zeta
 -2G_\sigma(\sigma,T)+\beta^{-1}G_\zeta(\sigma,T).
\]
The integral terms are \(O(\beta)\); all boundary terms are
exponentially small because \(u_R-\sigma\ge d_0>0\).
Thus \(J=O_{C^2_\sigma}(\beta)\). At the crossing,
\(\phi_\beta-\Gamma_\beta
=\beta\log(\omega_0/H_\beta(\sigma_\beta))=O_{C^2_W}(\beta)\).
Compose with the preceding \(C^2\) crossing map.
\end{proof}

\begin{lemma}\label{fx:lem-flat-part}
For some fixed \(d>0\) and finite integer \(n\),
\[
 \|\phi_\beta(u_+;W)-\phi_\beta(u_R;W)\|_{C^2}
          \le C\beta^{-n}e^{-d/\beta}.
\]
\end{lemma}
\begin{proof}
The strict cuts and preceding estimates give the fixed phase gap
\[
 \phi_\beta(u_R;W)\le\Gamma_{-,c}-3d_{\mathcal K}/4,\qquad
 \inf_{[u_R,u_+]}\Gamma_\beta\ge\Gamma_{-,c}-d_{\mathcal K}/2.
\]
The second bound follows from endpoint convergence and the fact that
\(\Gamma_\beta\) has at most one critical point, a maximum.
Equation~\eqref{fx:partial-fractions} gives bounded \(L/H_\beta\):
at the shrinking cut its only potentially large term satisfies
\(q/(a-u)\le q/(a\eta_+)\le C\). Thus \(F_\beta\) and its first
two \(r\)-derivatives are bounded; ordinary \(u\)-derivatives at
that cut are not needed.

Since \(h,K>0\), the shot continues while it stays in the rectangle.
The identity \(r=e^{(\phi_\beta-\Gamma_\beta)/\beta}\) and the phase
gap give \(r\le Ce^{-d/\beta}\): up to a first loss, the phase change
is at most \(C\int r\,du\), too small to close half the gap.
For its first two phase variations,
\(r_\phi=r/\beta\), \(r_{\phi\phi}=r/\beta^2\).
The scalar triangular variation argument therefore has exponentially
small integrated coefficients times finite powers of \(\beta^{-1}\).
It preserves the bounded \(W\)-variations at \(u_R\) and gives the
claimed \(C^2\) error.
\end{proof}

\begin{proposition}\label{fx:prop-inverse-phase}
The scalar backward-shooting solutions are defined up to \(u_+\),
and, jointly in the compact \(c\)-family, on every chord
\(\mathcal K_c\) above,
\begin{equation}\label{fx:phase-limit}
 \Psi_\theta\longrightarrow
       \Psi_{0,c}=\Gamma_c\circ\mathcal W^{-1}\quad\hbox{in }C^2 .
\end{equation}
More precisely, for the coefficient error \(\omega_{\rm d}\) in
Lemma~\ref{fx:lem-fast-inverse},
\begin{equation}\label{fx:phase-error-budget}
 \|\Psi_\theta-\Psi_{0,c}\|_{C^2(\mathcal K_c)}
 \le C\{\omega_{\rm d}+\beta(1+|\log\beta|)\}.
\end{equation}
Using \(\mathcal K_c=\mathcal J_c^\sharp\), its inverse
\(\mathcal H_\theta\) is defined on
\begin{equation}\label{fx:action-band}
 0\le\Theta\le\Gamma_{-,c}-\Delta
\end{equation}
and satisfies \(0<c_H\le\mathcal H_\theta'\le C_H\) and
\(|\mathcal H_\theta''/\mathcal H_\theta'|\le C_H\).
\end{proposition}
\begin{proof}
The exact derivative \eqref{fx:gamma-derivative}, together with
\(\Gamma_\beta(u_-)=\Gamma_c(u_-)=0\), gives
\(\Gamma_\beta-\Gamma_c=O_{C^3}(\omega_{\rm d}+\beta)\)
on the fixed negative interval.
Combining the three estimates in \eqref{fx:inverse-phase} and composing
with the uniformly \(C^2\) crossing \(\sigma_\beta(W)\) gives
\eqref{fx:phase-error-budget}; the exponentially small remainder
is absorbed into \(O(\beta)\). All derivatives here are in \(W\)
with the coefficients fixed.
For \(r=\mathcal W^{-1}(W)\),
\(\Psi_{0,c}'(W)=1/H_{0,c}(r)>0\) with uniform upper and lower
bounds. The larger intervals in \eqref{fx:W-limit} give a
fixed image margin around \eqref{fx:action-band}.
The uniform inverse function theorem applies. Explicitly,
\[
 u_{*,c}(\Theta)=
 \frac{c(1+u_-)e^\Theta-(c-u_-)}
      {(1+u_-)e^\Theta+(c-u_-)},\qquad
 \mathcal H_{0,c}(\Theta)=\frac{u_c^2-u_{*,c}(\Theta)^2}{2}.
\]
In particular, \(\Theta=0\) has a positive interior height on
the fixed exit section.
\end{proof}

The strict gap between the fixed-exit action endpoint and a negative
axis cut is
\begin{equation}\label{fx:overlap-gap}
 \Gamma_{-,c}-\Gamma_c(-\sigma)
    =\log\frac{c+\sigma}{c(1-\sigma)}>0,\qquad 0<\sigma<1 .
\end{equation}
Proposition~\ref{p2:prop-finite-v-fe-actual-budget} uses this gap to
place the fixed-exit construction and the bounded-height construction
on overlapping actual receivers. It then proves the completed-return
curvature estimate on their union, which is the physical interval used
by the coefficient cover.
